\chapter{Groundwater concepts}\label{ch:c2}\label{ch:concepts}
This chapter introduces the physical ideas and vocabulary used in the rest of the document. Readers familiar with
groundwater hydrology can skim it; everyone else will find every later equation built from the few ideas presented here.


\section{Head}
The \emph{head} $h$ at a point is the elevation to which water would rise in a small well drilled to that point, in
metres above a common datum (here, the same datum as the digital elevation model). Groundwater always flows from
higher head to lower head. In an unconfined aquifer the head is the elevation of the water table; the \emph{depth to
water} is land-surface elevation minus head.

\section{Darcy's law}
The flow of water through a porous material is proportional to the drop in head and to the cross-section, and
inversely proportional to the distance travelled:
\begin{equation}
Q = K\,A\,\frac{h_1-h_2}{L},
\end{equation}
where $Q$ is the flow [m$^3$/d], $K$ the hydraulic conductivity [m/d], $A$ the area the water passes through [m$^2$],
$h_1-h_2$ the head difference [m] and $L$ the distance [m].

\begin{example}[: Darcy's law]
Gravel with $K=50$ m/d, a 2\,000 m wide and 10 m thick strip ($A=20\,000$ m$^2$), heads 101 m and 100 m, 2\,000 m apart:
$Q=50\times20\,000\times1/2\,000=500$ m$^3$/d.
\end{example}

\section{Transmissivity and conductance}
Because $K$, $A$ and $L$ always appear together, it is convenient to combine them into a \emph{conductance}
$C = K A/L$ [m$^2$/d], so that $Q=C\,(h_1-h_2)$. Every exchange in this document --- between two grid cells, between an
aquifer and a river, a drain, a wetland --- is written in this form: \textbf{flow = conductance $\times$ head
difference}. For a layer of saturated thickness $b$, the \emph{transmissivity} is $T=Kb$ [m$^2$/d].

Between two neighbouring cells with different conductivities $K_1,K_2$, the conductance is the series combination
(harmonic mean), because water must pass through both halves in turn:
\begin{equation}
C_{12}=\frac{1}{\dfrac{L_1}{K_1A}+\dfrac{L_2}{K_2A}} .
\end{equation}
\begin{example}[: harmonic mean]
Two 2\,000 m cells, 10 m thick and 2\,000 m wide ($A=20\,000$ m$^2$, $L_1=L_2=1\,000$ m to the shared face), gravel
$K_1=50$ m/d next to till $K_2=0.2$ m/d:
$C_{12}=1/(1\,000/(50\cdot20\,000)+1\,000/(0.2\cdot20\,000)) = 1/(0.001+0.25)=3.98$ m$^2$/d. The low-conductivity
material controls the flow.
\end{example}

\section{Storage}
When the head in an aquifer falls, water is released from storage.
\begin{itemize}
\item In an \emph{unconfined} aquifer the water table drops and pores drain. The water released per square metre per
metre of fall is the \emph{specific yield} $S_y$ (typically 0.05--0.30).
\item In a \emph{confined} (full) aquifer, water is released only by the slight expansion of water and compression of
the grains; per cubic metre of aquifer per metre of head fall this is the \emph{specific storage} $S_s$
(typically $10^{-6}$--$10^{-4}$ 1/m).
\end{itemize}
\begin{example}[: storage change]
A 4 km$^2$ cell of gravel ($S_y=0.25$) whose water table falls 0.1 m in a day releases
$4\times10^6\times0.25\times0.1=100\,000$ m$^3$. The same fall in a 20 m thick confined sand ($S_s=10^{-5}$ 1/m)
releases $4\times10^6\times20\times10^{-5}\times0.1=80$ m$^3$ --- over a thousand times less.
\end{example}

\section{The groundwater flow equation}
Combining Darcy's law with conservation of water (what flows in, minus what flows out, equals the change in storage)
gives the equation MODFLOW solves. For one cell $n$ with neighbours $m$:
\begin{equation}
\sum_{m}C_{nm}\left(h_m-h_n\right) + \sum_{b}q_{b,n} = S_n\,\frac{h_n^{\,t+\dt}-h_n^{\,t}}{\dt},
\end{equation}
where all heads are those at the end of the step, $t+\dt$ (MODFLOW's implicit scheme, which is stable for any step
length), the first sum is flow from the neighbours, $q_{b,n}$ are the boundary flows of the cell (recharge, rivers,
drains, wells, \dots; positive into the aquifer), and $S_n$ is the storage capacity of the cell ($S_y$ times area when
unconfined, $S_s$ times volume when confined). Writing this for every cell gives a large set of equations for the
unknown heads, which MODFLOW solves every time step.

\begin{plain}
Each cell is a small water account: inflows from neighbours and boundaries, minus outflows, equals the change of water
stored. MODFLOW finds the heads that make every cell's account balance at the same time.
\end{plain}

\section{Head-dependent boundaries}
Many exchanges depend on the head: a river gives water to the aquifer when the river is higher than the water table,
and receives water when it is lower. MODFLOW writes such a boundary as
\begin{equation}
q_b = C_b\left(h_b-h_n\right),
\end{equation}
where $h_b$ is the boundary's water level. Because $q_b$ depends on the unknown head $h_n$, MODFLOW solves for the
boundary flow together with the heads; this keeps the solution stable even with large conductances. In MODFLOW's
internal notation each boundary contributes two numbers per cell, HCOF and RHS, and its flow is
\begin{equation}
q_b = \mathrm{HCOF}\cdot h_n-\mathrm{RHS}.
\end{equation}
The coupling uses exactly this formula, with MODFLOW's own HCOF and RHS, to compute every exchanged flow. It is the same
formula MODFLOW uses for its own budget, so the numbers exchanged with Raven are exactly the numbers MODFLOW accounts for.

\section{Water table and potentiometric surface}
In an unconfined aquifer the head is the water-table elevation and the pores above it are partly empty. In a confined
aquifer, sealed above by a low-permeability layer, water rises in a well above the top of the aquifer; the imaginary
surface joining these levels is the \emph{potentiometric surface}. The same aquifer can be confined in one place and
unconfined in another, or change over time as heads fall; MODFLOW's \emph{convertible} layers handle this switch
automatically by using $S_s$ while full and $S_y$ once the head falls below the top of the layer.

\begin{figure}[!htbp]\centering
\begin{tikzpicture}[x=1cm,y=1cm,lab/.style={font=\scriptsize,align=center},leg/.style={font=\scriptsize,anchor=west,align=left}]
\fill[yellow!25] (0,3) rectangle (9,4.35); \fill[gray!35] (0,2.3) rectangle (9,3);
\fill[orange!20] (0,0.6) rectangle (9,2.3); \fill[black!12] (0,0) rectangle (9,0.6);
\fill[white] (0,4.42) -- (9,4.28) -- (9,4.4) -- (0,4.5) -- cycle;
\draw[thick,brown!60!black] (0,4.42) -- (9,4.28);
\draw[dotted,red!70!black,very thick] (0,4.95) -- (9,4.5);
\draw[dashed,blue!70!black,thick] (0,3.95) -- (9,3.6);
\node[lab] at (4.3,3.3) {sand (unconfined)}; \node[lab] at (4.3,2.65) {clay (aquitard)};
\node[lab] at (4.3,1.45) {gravel (confined)}; \node[lab] at (4.3,0.3) {bedrock};
\draw[line width=2.5pt,black!70] (1.6,5.25) -- (1.6,1.8); \draw[line width=2.5pt,black!70,dash pattern=on 1.5pt off 1.2pt] (1.6,1.8) -- (1.6,1.0);
\fill[red!70!black] (1.6,4.87) -- ++(-0.13,0.16) -- ++(0.26,0) -- cycle;
\node[lab,anchor=south] at (1.6,5.3) {well A, screened in gravel};
\draw[line width=2.5pt,black!70] (6.8,5.25) -- (6.8,3.55); \draw[line width=2.5pt,black!70,dash pattern=on 1.5pt off 1.2pt] (6.8,3.55) -- (6.8,3.1);
\fill[blue!70!black] (6.8,3.69) -- ++(-0.13,0.16) -- ++(0.26,0) -- cycle;
\node[lab,anchor=south] at (6.8,5.3) {well B, screened in sand};
\draw[-{Stealth[length=2mm]},thick,black!60] (8.1,1.95) -- node[lab,left]{leakage} (8.1,3.2);
\draw[red!70!black] (9,4.5) -- (9.35,4.95) node[leg] {potentiometric surface\\ of the confined gravel};
\draw[brown!60!black] (9,4.28) -- (9.35,4.2) node[leg] {land surface};
\draw[blue!70!black] (9,3.6) -- (9.35,3.45) node[leg] {water table of the sand};
\end{tikzpicture}
\caption{The same place can have two different heads: the water table of the upper sand and the higher potentiometric surface of
the confined gravel below the clay. Water leaks upward through the clay because the gravel head is higher.}\label{fig:potsurf}
\end{figure}
\section{Three kinds of boundary conditions}
\begin{table}[!htbp]\centering
\caption{Three kinds of boundary conditions.}
\begin{tabular}{P{0.22\linewidth}P{0.368\linewidth}P{0.34\linewidth}}\toprule\hdr
Kind & Meaning & In this coupling\\\midrule
Specified flow & the flow is known (it does not depend on the head) & recharge, wells, reservoir seepage\\
Specified head & the head is known; the flow is whatever is needed & fixed-head boundaries (CHD)\\
Head-dependent flow & flow = conductance $\times$ (boundary level $-$ head) & rivers, seepage drains, wetlands, lakes, general-head boundaries, water-table ET\\\bottomrule
\end{tabular}
\end{table}
Head-dependent boundaries are the heart of a coupled model: they let the aquifer and the surface water ``negotiate'' the
exchange according to the water levels on both sides.

\section{The Newton method in simple terms}
The flow equations are non-linear: in an unconfined layer, transmissivity depends on the saturated thickness, which
depends on the unknown head. MODFLOW therefore solves them iteratively: guess heads, compute the imbalance of every cell,
correct the heads, and repeat until the corrections are smaller than a tolerance (here 0.00001 m, i.e.\ 0.01 mm). The Newton--Raphson method
computes each correction from the slope of the imbalance with respect to the heads, which converges quickly even when cells
are nearly dry. Each repetition is an \emph{outer iteration}; \cmd{GWBudget.csv} reports how many were needed each step
(typically 3--10 in the Liard example).

\section{A cell's water account, with numbers}
\begin{example}[: one cell for one day]
A 4 km$^2$ gravel cell receives 8\,000 m$^3$ recharge; it gets 1\,500 m$^3$ from an up-gradient neighbour and passes
2\,500 m$^3$ to a down-gradient one; its drain removes 4\,000 m$^3$. Inflow $-$ outflow $=8\,000+1\,500-2\,500-4\,000=3\,000$
m$^3$ are stored, so with $S_y=0.25$ the water table rises by $3\,000/(4\times10^6\times0.25)=3$ mm. MODFLOW finds heads for
which every cell's account balances like this one at the same time.
\end{example}

\section{Units and conversions}
\begin{table}[!htbp]\centering
\caption{Units and conversions.}
\begin{tabular}{P{0.44\linewidth}P{0.5\linewidth}}\toprule\hdr
Quantity & Conversion\\\midrule
depth of water in an HRU to volume & $V\,[\mathrm{m}^3]=x\,[\mathrm{mm}]\times A\,[\mathrm{km}^2]\times1000$\\
flow & 1 m$^3$/s $=86\,400$ m$^3$/d\\
recharge rate & 1 mm/d over 1 km$^2$ $=1\,000$ m$^3$/d\\
hydraulic conductivity & 1 m/d $\approx1.16\times10^{-5}$ m/s\\
leakance & conductance per area, d$^{-1}$; $C=L\times$area\\\bottomrule
\end{tabular}
\end{table}
